% ============================================================
% §VI-C draft — accounting reframing. Written against
% results/tables_deployment_audit.tex (tab:dep:systems, tab:dep:sensitivity).
% Forward pointers: §VI-E (envelope sweep), §VI-F (early-exit mechanism).
% Every number here is traceable to results/deployment/*.csv.
% ============================================================

\subsection{Operation-Priced Energy Does Not Survive Measurement}
\label{sec:vi:accounting}

The efficiency claim in Section~\ref{sec:energy}, and the comparable claims that
are standard in the spiking-IDS literature, rest on an operation-priced proxy:
synaptic operations are counted and multiplied by a published per-operation
energy. We adopt that accounting throughout this paper. Having measured the
deployed systems directly, we now report that the proxy does not support the
conclusion we drew from it, for three separate reasons. We state them in the
order of how much of the discrepancy each explains.

\subsubsection{Priced consistently, the tree ensemble is not more expensive}
The proxy has been applied asymmetrically. A tree comparison is a
subtract-and-branch and belongs to the same Horowitz cost class as an
accumulate, yet tree-based baselines are conventionally excluded from
operation-priced comparison on the grounds that their operations are not
multiply--accumulates. Pricing every system on the same scale
(Table~\ref{tab:dep:systems}) removes that asymmetry, and it reverses the
ordering: the compact ensemble requires $2{,}768$\,pJ of arithmetic per sample
against $5{,}654$\,pJ for the deployed spiking student, a factor of $0.49$. The
full teacher, at $6{,}873$\,pJ, is within $1.22\times$ of the student rather than
orders of magnitude above it. The spiking model's arithmetic advantage is real
only against the \emph{dense} baselines, where the $4.6$ versus $0.9$\,pJ
differential between a multiply--accumulate and an accumulate applies.

\subsubsection{Arithmetic is a small fraction of the energy that is actually spent}
Extending the model to include memory traffic --- weight or node reads, membrane
state, and input reads, priced by a tier keyed to artifact size --- places
between $95.7\%$ and $99.9\%$ of total energy in the memory term for every system
measured. An accounting that prices only arithmetic therefore prices under $5\%$
of the system. This is the dominant source of error, and it is not specific to
spiking models: it applies identically to the dense and tree baselines, and it is
why the ensemble's real disadvantage is its $5.04$\,MB of node arrays rather than
its operation count.

\subsubsection{The remaining ordering depends on a modelling choice, not on measurement}
Once memory dominates, the comparison between the spiking student and its
parameter-identical dense twin is decided by the assumed reuse behaviour rather
than by either system's operation count. Table~\ref{tab:dep:sensitivity} reports
the ratio under two assumptions. Charging one cold access per operation, the
ordering \emph{inverts} above roughly $20$\,pJ per access and the dense twin is
cheaper. Charging the first touch of a distinct item at the tier price and
subsequent touches at an L1 cost, the spiking student is cheaper throughout the
swept range. Both are defensible: the spiking student re-reads a $53$\,KB weight
set across timesteps and is plausibly resident, whereas the ensemble walks
$76{,}475$ nodes through a multi-megabyte array with data-dependent branching and
is plausibly not. We report the band, $0.94$--$2.06\times$, and note explicitly
that it spans a modelling choice rather than measurement noise. We do not select
the assumption that favours our system.

\subsubsection{The per-operation constant is an arithmetic-only idealisation}
The $0.9$\,pJ figure used here and throughout the spiking literature prices
arithmetic alone. Measured neuromorphic silicon carries routing, spike delivery
and core-idle overheads roughly $25\times$ larger per synaptic operation. We have
not measured either system on such a device, and the comparison that would be
required --- a dense network on the same substrate --- does not exist, so we draw
no conclusion about relative performance there. The narrower and supportable
statement is that any deployment claim must be validated against the target's own
per-operation figure rather than against an analytical floor.

\subsubsection{What we therefore claim}
We withdraw the claim that Tree-to-Spike is the more energy-efficient choice on
general-purpose hardware. On a conventional processor with adequate memory the
gradient-boosted ensemble is both more accurate and lower-latency
(Table~\ref{tab:dep:systems}), and the dense twin is lower-latency by an order of
magnitude. What survives measurement is a memory-envelope argument, whose scope
is established in Section~\ref{sec:vi:envelope}, and it is conditional: it holds
where the taxonomy is large enough or hard enough that an ensemble cannot
converge within the envelope. The efficiency claim is thus not merely reduced in
magnitude but changed in kind, from an operation-count advantage to a footprint
constraint.

We regard the methodological point as the more durable contribution. An
analytical proxy can fail in two distinct ways, and they are distinguishable by
test. It may be \emph{incomplete}, omitting a term that dominates --- detected by
adding the omitted term and observing the estimate move, as in
Section~\ref{sec:vi:accounting}\ref{subsec:memory}. Or it may be
\emph{target-mismatched}, faithfully estimating a different quantity --- detected
by exhausting the statistic space over the same information and observing that
nothing moves, as in the early-exit analysis of
Section~\ref{sec:vi:earlyexit}. The two standard proxies of this literature,
operation-priced energy and prefix confidence, exemplify one failure mode each.
